\section{Exact cyclotomic base change and the different}\label{sec:different}
This section proves Theorems~\ref{thm:cyclotomic-exact} and~\ref{thm:cyclotomic-bound} for every odd prime.

For a finite totally ramified extension $A/B$, write
$d(A/B)=v_A(\mathfrak D_{A/B})$. We use
\begin{equation}\label{eq:different-tower}
 d(A/C)=d(A/B)+[A:B]d(B/C).
\end{equation}
This is \citet[Chapter~III, \S4, Proposition~8, p.~51]{Serre1979}. When $\Ocal_A=\Ocal_B[\pi]$ and $f$ is the minimal polynomial of $\pi$, the different is generated by $f'(\pi)$ \citep[Chapter~III, \S6, Proposition~11 and Corollary~2, pp.~55--57]{Serre1979}. For a finite Galois extension with lower ramification groups $G_i$,
\begin{equation}\label{eq:different-lower}
 d(A/B)=\sum_{i\ge0}(|G_i|-1),
\end{equation}
as in Chapter~IV, \S1, Proposition~4, p.~64 of that source. These statements require no finite-residue-field hypothesis.

\begin{proof}[Proof of Theorem~\ref{thm:cyclotomic-exact}]
A Cohen coefficient field gives $[K:K_0]=p(p-1)$, hence $[K:K_1]=p$. The Eisenstein derivative calculation given below proves $d\le p^2-1$. Let $T/K_1$ be the normal closure of $K/K_1$. Its Galois group acts faithfully and transitively on the $p$ roots of a minimal polynomial. Since the residue field is algebraically closed, the entire group is inertia. Its wild subgroup is a nontrivial normal $p$-subgroup of the symmetric group on $p$ letters, hence has order $p$ and acts regularly. The normalizer of this subgroup is the affine group $C_p\rtimes\F_p^\times$. Thus
\[
 \Gal(T/K_1)=C_p\rtimes C_m,\qquad m\mid p-1,
\]
with faithful tame action when $m>1$. The extension $T/K$ is tame of degree $m$. If $b$ is the unique positive lower break of $T/K_1$, the different formula in the tower $T/K/K_1$ gives
\[
 pm-1+(p-1)b=(m-1)+md,
 \qquad u_T=\frac bm=\frac{d-(p-1)}{p-1}\le p.
\]
Here $u_T$ is its unique positive upper break.

Set $C_0=K_0(\mu_{p^3})$. The cyclic extension $C_0/K_1$ has upper breaks $p-1$ and $2(p-1)$: its lower breaks, computed on $\zeta_{p^3}-1$, are $p-1$ and $p^2-1$. It is disjoint from $T$ over $K_1$. Indeed, a nontrivial intersection would contain $K_0(\mu_{p^2})$ and give a cyclic quotient of order $p$ of $C_p\rtimes C_m$. If $m>1$, the faithful tame action makes the $C_p$ factor trivial in the abelianization. If $m=1$, the intersection would force $K=K_0(\mu_{p^2})$, contrary to the hypothesis.

Since $2(p-1)>u_T$, the last cyclotomic ramification group persists after the $T$ factor has become trivial. The quotient and subgroup forms of Herbrand's theorem therefore give the last upper break of $TC_0/T$ as
\begin{align*}
 \psi_{T/K_1}\bigl(2(p-1)\bigr)
 &=m u_T+pm\bigl(2(p-1)-u_T\bigr)\\
 &=m\bigl((p-1)(2p+1)-d\bigr).
\end{align*}
The tame extension $T/K$ has degree $m$ and is disjoint from the cyclotomic extension over $K$. Tame base change multiplies its positive upper breaks by $m$. Dividing the displayed expression by $m$ proves the identity; $d\le p^2-1$ gives the lower bound $p(p-1)$.
\end{proof}

\begin{proof}[Proof of Theorem~\ref{thm:cyclotomic-bound}]
Choose a Cohen coefficient ring $W(k)\subset\Ocal_K$ and write $K_0=W(k)[1/p]$. The extension $K/K_0$ is finite of degree $e_K$. For completeness, if $\pi$ is a uniformizer of $K$, then $1,\pi,\ldots,\pi^{e_K-1}$ give a basis of $\Ocal_K/p\Ocal_K$ over $k$. Successive $p$-adic lifting spans $\Ocal_K$ over $W(k)$, and distinct valuations modulo $e_K$ prove independence. This argument works for infinite perfect $k$.

Put
\[
 K_1=K_0(\mu_p),\qquad C_0=K_0(\mu_{p^3}),\qquad C=KC_0.
\]
Cyclotomic Eisenstein polynomials give $[K_1:K_0]=p-1$ and $[K:K_1]=p$. For odd $p$, $C_0/K_1$ is cyclic of degree $p^2$ and its unique degree-$p$ subfield is $K_0(\mu_{p^2})$. The hypothesis $\mu_{p^2}\not\subset K$ therefore implies $K\cap C_0=K_1$. Consequently
\[
 [C:K]=p^2,\qquad[C:C_0]=p.
\]
All fields have the same residue field and all extensions in these towers are totally ramified.

We need two elementary bounds:
\begin{equation}\label{eq:different-bounds}
 d(K/K_1)\le p^2-1,
 \qquad d(C/C_0)\ge p-1.
\end{equation}
A uniformizer of $K$ generates its ring of integers over $\Ocal_{K_1}$. Its minimal polynomial is Eisenstein,
$f(X)=X^p+\sum_{i=0}^{p-1}a_iX^i$.
The leading derivative term $p\pi^{p-1}$ has valuation $p(p-1)+(p-1)=p^2-1$. Each other nonzero term $ia_i\pi^{i-1}$ has valuation congruent to $i-1$ modulo $p$, because $1\le i<p$ and $v_K(a_i)$ is divisible by $p$. These residues are pairwise distinct and exclude the residue $-1$ of the leading term. Thus the least valuation occurs only once and cannot cancel, proving the first inequality. For the degree-$p$ extension $C/C_0$, the Eisenstein derivative formula gives valuation at least $p-1$ for every term, proving the second.

The cyclotomic extension $C_0/K_1$ has lower breaks $p-1$ and $p^2-1$. Indeed, $\zeta_{p^3}-1$ is a uniformizer. An exponent $1+pa$, with $p\nmid a$, displaces it by an element of valuation $p$; an exponent $1+p^2a$, with $p\nmid a$, gives valuation $p^2$. Summing these displacements gives
\[
 d(C_0/K_1)=(p^2-p)p+(p-1)p^2=2p^2(p-1).
\]
Different transitivity in the two towers over $K_1$ yields
\begin{align*}
 d(C/K)+p^2d(K/K_1)
 &=d(C/C_0)+p\,d(C_0/K_1),\\
 d(C/K)&\ge p-1+2p^3(p-1)-p^2(p^2-1)\\
 &=p^2(p-1)^2+p-1.
\end{align*}

A totally ramified cyclic group of order $p^2$ has two positive breaks. In terms of its upper breaks $\delta_1<\delta_2$, the lower breaks are $\delta_1$ and $p\delta_2-(p-1)\delta_1$. The upper-to-lower conversion is the Herbrand function of \citet[Chapter~IV, \S3, pp.~73--76]{Serre1979}. Summing the lower groups gives
\[
 d(C/K)=(p^2-1)(\delta_1+1)
 +(p-1)\bigl(p\delta_2-p\delta_1\bigr)
 =(p-1)(\delta_1+p\delta_2+p+1).
\]
Substitution of the preceding lower bound gives
$\delta_1+p\delta_2\ge p(p^2-p-1)$. If $\delta_1\le p-1$, divide the resulting inequality for $p\delta_2$ by $p$ to obtain~\eqref{eq:second-break}.
\end{proof}

\begin{remark}
No integrality theorem for upper breaks was used. The two-tower estimate alone gives $\delta_2\ge91/5$ at $p=5$ when $\delta_1\le4$. The exact normal-closure argument gives the unconditional stronger bound $\delta_2\ge20$ under the stated hypotheses. Neither argument needs Hasse--Arf.
\end{remark}
